% TeX DSL 国际化方案
% 
% 核心理念：
% 1. 编剧直接写原文，不需要关心 ID
% 2. 翻译工具自动提取文本
% 3. 翻译团队只翻译提取的文本文件

% ============================================
% 编剧直接写原文
% ============================================

\scene[school_gate.png, theme.mp3]

星光学院，这所历史悠久的名校，隐藏着不为人知的秘密。

我转学到这里已经一周了，却总觉得有什么不对劲。

\flash

\p[thinking] （奇怪...那个旧校舍，为什么总是锁着？）

\role[sakura][curious, left]

\s 新同学！你在看什么呢？

\p[surprised] 啊，樱井同学...没什么，只是在发呆。

\s[smile] 你的表情可不像在发呆哦。对了，放学后学生会有个活动，要来吗？

```choice
\case[accept][+s] 好啊，我去看看
\case[decline][-s] 抱歉，我还有事
\case[ask] 学生会有什么活动？
```

% ============================================
% 节点 ID 设计（替代行号）
% ============================================

% 为什么不用行号？
% 行号在剧本修改后会变化，导致翻译引用失效。

% ID 格式：<文件路径>:<结构路径>#<内容指纹>
%
% 示例：
%   main.tex:scene[0].narration[0]#a3b2c1d0
%   main.tex:scene[0].dialogue[2]#e4f5g6h7
%   main.tex:scene[0].choice[0].option[1]#i8j9k0l1

% 组成部分：
% - 文件路径：源文件相对路径
% - 结构路径：节点在 AST 中的位置
% - 内容指纹：文本内容的 SHA-256 哈希前 8 位

% ============================================
% 翻译工作流
% ============================================

% 步骤 1：编剧用中文写完剧本
% 步骤 2：运行提取工具
%   $ galgame extract main.tex -o locales/zh-CN/main.pot
%
% 步骤 3：生成各语言翻译文件
%   $ galgame init-locale zh-CN main.pot  # 源语言，直接复制
%   $ galgame init-locale en-US main.pot  # 英文翻译文件
%   $ galgame init-locale ja-JP main.pot  # 日文翻译文件
%
% 步骤 4：翻译团队编辑翻译文件
% 步骤 5：运行时根据语言设置加载翻译

% ============================================
% 生成的翻译文件格式
% ============================================

# locales/en-US/main.po (gettext 格式)

#. ID: main.tex:scene[0].narration[0]#a3b2c1d0
msgid "星光学院，这所历史悠久的名校，隐藏着不为人知的秘密。"
msgstr "Star Academy, a historic school, hides secrets unknown to all."

#. ID: main.tex:scene[0].narration[1]#b4c5d6e7
msgid "我转学到这里已经一周了，却总觉得有什么不对劲。"
msgstr "I transferred here a week ago, but something feels off."

#. ID: main.tex:scene[0].dialogue[0]#c5d6e7f8
msgctxt "player"
msgid "（奇怪...那个旧校舍，为什么总是锁着？）"
msgstr "(Strange... why is that old building always locked?)"

#. ID: main.tex:scene[0].dialogue[1]#d6e7f8g9
msgctxt "sakura"
msgid "新同学！你在看什么呢？"
msgstr "New student! What are you looking at?"

#. ID: main.tex:scene[0].dialogue[2]#e7f8g9h0
msgctxt "sakura"
msgid "你的表情可不像在发呆哦。对了，放学后学生会有个活动，要来吗？"
msgstr "You don't look like you're spacing out. By the way, there's a student council activity after school. Want to come?"

# 选择选项
#. ID: main.tex:scene[0].choice[0].option[0]#f8g9h0i1
msgid "好啊，我去看看"
msgstr "Sure, I'll check it out"

#. ID: main.tex:scene[0].choice[0].option[1]#g9h0i1j2
msgid "抱歉，我还有事"
msgstr "Sorry, I have something else"

#. ID: main.tex:scene[0].choice[0].option[2]#h0i1j2k3
msgid "学生会有什么活动？"
msgstr "What kind of activity?"

% ============================================
% 增量更新策略
% ============================================

% 当剧本修改后，使用智能匹配保留已有翻译：

% 1. 精确匹配：ID 完全相同 → 直接保留翻译
% 2. 指纹匹配：指纹相同，路径变化 → 更新路径，保留翻译
% 3. 模糊匹配：路径相似，文本相似度 > 80% → 标记为 fuzzy
% 4. 新增条目：无法匹配 → msgstr 为空，等待翻译
% 5. 废弃条目：旧 ID 不在新提取中 → 标记为 obsolete

% ============================================
% 场景演示：插入内容后更新
% ============================================

% 原始剧本：
% \scene[school.png]
% 星光学院，历史悠久的名校。
% \s 你好！

% 翻译文件：
% #. ID: main.tex:scene[0].narration[0]#abc12345
% msgid "星光学院，历史悠久的名校。"
% msgstr "Star Academy, a historic school."
%
% #. ID: main.tex:scene[0].dialogue[0]#def67890
% msgctxt "sakura"
% msgid "你好！"
% msgstr "Hello!"

% 修改后剧本（插入内容）：
% \scene[school.png]
% 阳光明媚的早晨。
% 星光学院，历史悠久的名校。
% \s 你好！

% 增量更新结果：
% #. ID: main.tex:scene[0].narration[0]#xyz11111  # 新增
% msgid "阳光明媚的早晨。"
% msgstr ""  # 等待翻译
%
% #. ID: main.tex:scene[0].narration[1]#abc12345  # 路径更新，指纹匹配
% msgid "星光学院，历史悠久的名校。"
% msgstr "Star Academy, a historic school."  # 保留翻译
%
% #. ID: main.tex:scene[0].dialogue[0]#def67890  # 精确匹配
% msgctxt "sakura"
% msgid "你好！"
% msgstr "Hello!"  # 保留翻译

% ============================================
% JSON 格式的翻译文件（备选）
% ============================================

# locales/en-US/main.json

{
  "meta": {
    "source": "main.tex",
    "language": "en-US",
    "version": "1.0.0"
  },
  "translations": [
    {
      "id": "scene[0].narration[0]#a3b2c1d0",
      "context": "narration",
      "source": "星光学院，这所历史悠久的名校，隐藏着不为人知的秘密。",
      "target": "Star Academy, a historic school, hides secrets unknown to all."
    },
    {
      "id": "scene[0].narration[1]#b4c5d6e7",
      "context": "narration",
      "source": "我转学到这里已经一周了，却总觉得有什么不对劲。",
      "target": "I transferred here a week ago, but something feels off."
    },
    {
      "id": "scene[0].dialogue[0]#c5d6e7f8",
      "context": "player.thinking",
      "source": "（奇怪...那个旧校舍，为什么总是锁着？）",
      "target": "(Strange... why is that old building always locked?)"
    },
    {
      "id": "scene[0].dialogue[1]#d6e7f8g9",
      "context": "sakura",
      "source": "新同学！你在看什么呢？",
      "target": "New student! What are you looking at?"
    }
  ]
}

% ============================================
% 角色名称翻译
% ============================================

# locales/en-US/characters.json

{
  "player": "Protagonist",
  "sakura": "Sakura Mirai",
  "ren": "Ren Kurosaki",
  "yuki": "Yuki Shirayuki"
}

# locales/ja-JP/characters.json

{
  "player": "主人公",
  "sakura": "桜井美咲",
  "ren": "黒崎蓮",
  "yuki": "白雪悠希"
}

% ============================================
% 工作流对比
% ============================================

❌ 错误方式（使用行号）：
#: main.tex:15        # 在第 15 行插入内容后，这里变成 16 行
msgid "星光学院..."
msgstr "..."

✅ 正确方式（使用节点 ID）：
#. ID: main.tex:scene[0].narration[0]#a3b2c1d0
msgid "星光学院..."
msgstr "..."

% ============================================
% 完整工作流程
% ============================================

1. 编剧阶段
   - 编剧用中文（或任何源语言）直接写剧本
   - 不需要关心翻译，只关注故事内容
   - 文件：scripts/main.tex

2. 提取阶段
   - 运行：galgame extract scripts/*.tex -o locales/zh-CN/
   - 自动提取所有可翻译文本
   - 生成带节点 ID 的翻译文件

3. 翻译阶段
   - 翻译团队编辑 PO/JSON 文件
   - 不需要触碰剧本文件
   - 文件：locales/en-US/main.po, locales/ja-JP/main.po

4. 构建阶段
   - 运行：galgame build --locale en-US
   - 将翻译合并到运行时数据

5. 运行时
   - 根据用户语言设置显示对应翻译
   - 没有翻译的文本显示原文
